\begin{informalproof}
The step width $s$ bounds the joint record at a step, including the state and the work symbols needed to continue
the run. Write $n(d)\ge1$ for an integer upper bound on the outputs per component; replacing it by a larger integer
does not change $n(d)=o(\log\log d)$. Choose constants $0<\delta<1$ and $\eta>0$ such that
\[
 (k+2)(s\theta c+\eta)<(1-\delta)/2.
\]
In addition to the assumed cut $\Gamma_h$, use the depth reduction for oblivious machines~\cite[Corollary~9.3]{BoundedStack2026} to obtain a cut
$\Gamma_D$ of record cost at most $\eta L$ and a remaining depth $D$ independent of $d$. For
$\Gamma=\Gamma_h\cup\Gamma_D$ the undirected components only shrink, both conclusions hold, and
\begin{equation}\label{eq:fragment-upper}
 s\card\Gamma\le s\theta c(L+1)+\eta L.
\end{equation}
Because the schedule is oblivious these cuts and the ensuing helper sets can be fixed before the input.

Each output now has at most $B=\sum_{j=0}^{D}(k+1)^j$ raw input helpers. Apply the paid degree reduction of
\cite[Corollary~3.8]{DirectedRecovery2026}, conditioning on input coordinates that occur in too many cones.
For a fixed sufficiently large reverse-degree cutoff this retains at least $(1-\delta)L$ demanded coordinates,
and their paired-port incidence relation has both degrees at most a constant $\Delta\ge2$. The native
paired-distance estimate of \cite[Theorem~3.3]{DirectedRecovery2026} is needed with a slowly growing radius here,
not merely in its fixed-radius asymptotic form. Its finite packing proof gives, at radius $2n-1$,
\[
 \#\{\text{retained targets at paired distance at most }2n-1\}
 \ \le\ V_{2n-1}\!\sum_{j=1}^{n}2j\,\bigl(a(2j-1)+3\bigr)
              \frac{L}{\log_2d},
\]
where $V_r=\sum_{j=0}^{r}\Delta^j$ and $a=4c_*(B+1)(k+2)(k+6)$. Here $c_*=2c$ ensures
$T\le c_*L$ for $L\ge1$, as required by that finite estimate. The coefficient is
$O_{k,D,c,\Delta}(n^3\Delta^{2n})$. Since $n=o(\log\log d)$, its ratio to $\log_2d$ tends to zero. Remove these
$o(L)$ targets, leaving at least $(1-\delta)L-o(L)$ good ones. This step uses the explicit transport estimate of the fifth paper~\cite{DirectedRecovery2026}, whose constants are
proved there in the text and are not formalized; only its finite packing step is in Lean.

Partition the good outputs according to the uncut component containing their output step, with a singleton patch
for a good output at a cut root. A patch has at most $n$ targets. Its induced paired-port graph has at most $2n$
vertices. If the two ports of a target were connected there, a shortest path would have length at most $2n-1$,
contrary to its being good. Thus the separated-components lemma~\cite{DirectedRecovery2026} bounds every storage
code on the patch's induced relation by $2^{|J|/2}$.

For each patch fix all exterior input bits, including the paid conditioning coordinates. Let $T_J$ be the cut
records at its incoming boundary; include the root record for a cut-root output. Local determination on the DAG
makes these records an index code for the demands in $J$, with the induced raw helpers. The fibre bound gives
$s|T_J|\ge |J|/2$. The records need not be independent and are not charged as a joint entropy. This is precisely
the abstract local-view invoice, formalized as \lean{LocalRecordViews.half_invoice}.

A cut vertex has at most $k+1$ forward children, each entering at most one uncut component, and at most one output
of its own. Hence $\sum_J|T_J|\le(k+2)|\Gamma|$. Summing the patch inequalities gives
\[
 s(k+2)|\Gamma|\ge(1-\delta)L/2-o(L),
\]
contradicting \eqref{eq:fragment-upper} for large $d$. The argument uses outputs per component, not its number of
steps or its time span. No implication in either direction with the macroscopic-motion classes of
\cite{FourthPaper2026} is established.
\end{informalproof}
